\scanpage{435}

\unsure{linked} in a much more natural way to the very structure
of the group $\mathcal{M}_{\rho,\sigma}$ and to its homomorphisms
$\xrightarrow{\Theta}G$ and $\chi=\det\circ\Theta$ into $G_m$, namely the
\unsure{integral} subgroup $SR_{\rho,\sigma}$, allowing us to write
even $\mathcal{M}_{\rho,\sigma}^{+}$ as a
\uline{direct product}
\[
\tag{99}
\begin{array}{rccccc}
\mathcal{M}_{\rho,\sigma}^{+}&\simeq&SR_{\rho,\sigma}&\times&
 \unsure{\mathcal{G}|\mathcal{C}}_{\rho,\sigma}^{+}\\
&&\rotatebox{90}{$\simeq$}&&\rotatebox{90}{$\unsure{=|\simeq}$}\\
&&S'T_\rho\times SN'_\sigma&&SG
\end{array}
.
\]
inducing an analogous decomposition
\[
\tag{100}
\left\{
\begin{aligned}
\mathcal{M}_{\rho,\sigma}^{+\circ}
 &\simeq \underset{ST_\rho\times ST_\sigma}{\underbrace{
 SR_{\rho,\sigma}^{\deleted{\unsure{+}}\inserted{\circ}}}}
 \times
 \underset{SG}{\underbrace{\deleted{\unsure{S}}\unsure{\mathcal{G}|\mathcal{C}}_{\rho,\sigma}^{+}}}\\
 &\simeq ST_\rho\times ST_\sigma\times SG,
\end{aligned}
\right.
\]
\margin{this is (46), which I had forgotten along the way\ldots}
who can do better!. On the other hand, the fact that
$R_{\rho,\sigma}\longrightarrow\mathcal{T}_{\rho,\sigma}$ is an
\uline{isogeny} mod $\mu_2$ shows that we have, up to this isogeny, a
``section'' \unsure{of|even}
\deleted{of $\unsure{\mathcal{T}_\rho}$}
\inserted{$\mathcal{M}_{\rho,\sigma}$ itself}
as an extension of $\mathcal{T}_{\rho,\sigma}$ by
$\unsure{\mathcal{G}|\mathcal{C}}_{\rho,\sigma}^{+}\simeq SG$.
In fact, we have
\[
\tag{101}
\left\{
\begin{array}{l}
R_{\rho,\sigma}\cap SG=\mu_2\\[4pt]
\deleted{\text{\unsure{but} }R_{\rho,\sigma}\text{ \unsure{no longer commutes},}}\\[-2pt]
\deleted{\text{\unsure{so we have an exact sequence}}}\\[4pt]
R_{\rho,\sigma}\text{ commutes with }SG
\end{array}
\right.
\]
\hfill (trivial by definition of $R_{\rho,\sigma}$ as
$\operatorname{Ker}\widetilde{\Theta}$

so we have an exact sequence

\mspage{434}

\begin{center}
{\small
434\\
\texttt{https://grothendieck.umontpellier.fr}\\
Université de Montpellier, UMR CNRS 5149, Case courrier 051, Place Eugène Bataillon\\
34095 Montpellier cedex 5 - France
}
\end{center}

\scanpage{436}
\mspage{520}

\noindent(102)
\diagram{\begin{tikzcd}[ampersand replacement=\&,column sep=large,row sep=small]
1\arrow[r]\&\mu_2\arrow[r]\&R_{\rho,\sigma}\times SG\arrow[r]\&
 \mathcal{M}_{\rho,\sigma}\arrow[r]\&1\\
\&\varepsilon\arrow[r,mapsto]\&
 \left(\varepsilon,
 \underset{\inserted{\substack{\Vert\\\varepsilon^{-1}}}}
 {\varepsilon^{\deleted{\unsure{+}}}}\right)
\end{tikzcd}}

and likewise

\noindent(103)
\diagram{\begin{tikzcd}[ampersand replacement=\&,column sep=large]
1\arrow[r]\&\mu_2\arrow[r]\&R_{\rho,\sigma}^{\circ}\times SG\arrow[r]\&
\mathcal{M}_{\rho,\sigma}^{\circ}\arrow[r]\&1
\end{tikzcd}}

already considered (a little timidly) in (78).
It is decidedly time to make fully explicit the description of
$R_{\rho,\sigma}$ in terms of $T_\rho,N'_\sigma\ldots$ --- we obtain it,
in the context of page~510, \deleted{as} as
\[
\tag{104}
\begin{aligned}
R_{\rho,\sigma}
={}&\left\{
 \left(
 \underset{\substack{\rotatebox{90}{$\in$}\\\Gamma T_0}}{\underline u},
 \underset{\substack{\rotatebox{90}{$\in$}\\\Gamma T_\rho}}{\underline a},
 \underset{\substack{\rotatebox{90}{$\in$}\\\Gamma N'_\sigma}}{\underline b}
 \right)\cdot
 \underset{\substack{\rotatebox{90}{$\in$}\\SG}}{f}
 \ \middle|\ 
 \begin{gathered}
 \det\underline u=\det\underline b=\varepsilon'\det\underline a,
 \quad\varepsilon'{}^2=1\\
 \underline u f=\lambda\qquad(\varepsilon',\lambda\in\Gamma G_m)\\
 (\text{hence }\lambda^2=\det\underline u)
 \end{gathered}\right\}\\[8pt]
\simeq{}&\left\{
 \left(\underset{T_0}{\underline u},
       \underset{T_\rho}{\underline a},
       \underset{N'_\sigma}{\underline b}\deleted{\unsure{)}},
       \underset{G_m}{\lambda}\right)
 \ \middle|\ \det\underline u=\det\underline b
 =\varepsilon'\det\underline a=\lambda^2\right\}\\[8pt]
\simeq{}&\left\{
 \left(\underset{T_\rho}{\underline a},
       \underset{N'_\sigma}{\underline b},
       \underset{G_m}{\lambda}\deleted{\unsure{)}},
       \inserted{\underset{\mu_2}{\varepsilon'}}\right)
 \ \middle|\
 \begin{gathered}
 \inserted{\det\underline b=\varepsilon'\det\underline a=\lambda^2}\\
 \deleted{\unsure{\det\underline b/\det\underline a\in\Gamma\mu_2}}
 \end{gathered}\right\}\\[8pt]
\simeq{}&\left\{
 \left(\underset{T_\rho}{\underline a},
       \underset{N'_\sigma}{\underline b},
       \underset{G_m}{\lambda}\right)
 \ \middle|\ \det\underline b=\lambda^2,\quad
 \det\underline a^2=\det\underline b^2\right\}\\[8pt]
\simeq{}&\left(\deleted{\unsure{\{}}(T_\rho,\delta_\rho^2)
 \times_{G_m}(N'_\sigma,\delta_\sigma^2)\right)_{(\delta_\sigma)}
 \times_{G_m}G_m{}_{(\delta^2)}
\end{aligned}
\]
or more simply, \deleted{which gives}
\inserted{in terms of $\mathcal{T}_{\rho,\sigma}$:}
\[
\deleted{(105)\qquad
\begin{aligned}
 R_{\rho,\sigma}^{\circ}
 &\simeq\left\{
 \left(\underset{T_\rho}{\underline a},
       \underset{T_\sigma}{\underline b},
       \underset{G_m}{\lambda}\right)
 \ \middle|\ \det\underline a=\det\underline b=\lambda^2
 \right\}\\
 &\simeq\unsure{\mathcal{G}|\mathcal{C}}_{\rho,\sigma}^{\circ}
\end{aligned}}
\]

\noindent(105)
\diagram{\begin{tikzcd}[ampersand replacement=\&,column sep=large,row sep=large]
R_{\rho,\sigma}\arrow[r]\arrow[d,"{\kappa}"']\&
 \mathcal{T}_{\rho,\sigma}\arrow[d,"{\chi}"]\&\text{diagram}\\
G_m\arrow[r,"{\delta^2}"']\&G_m\&\text{cartesian}\\[-1em]
\lambda\arrow[r,mapsto]\&\lambda^2
\end{tikzcd}},

\mspage{435}

\begin{center}
{\small
435\\
\texttt{https://grothendieck.umontpellier.fr}\\
Université de Montpellier, UMR CNRS 5149, Case courrier 051, Place Eugène Bataillon\\
34095 Montpellier cedex 5 - France
}
\end{center}

\scanpage{437}

fitting into a diagram of exact sequences

\noindent(106)
\diagram{\begin{tikzcd}[ampersand replacement=\&,column sep=large,row sep=large]
\&\&1\arrow[d]\&1\arrow[d]\&\\
\deleted{\unsure{1}}\arrow[r,draw=none,"{\deleted{\longrightarrow}}" description]\&
\deleted{\unsure{\mu_2}}\arrow[r,draw=none,"{\deleted{\longrightarrow}}" description]\&
SR_{\rho,\sigma}\arrow[r,"{\simeq}"]\arrow[d]\arrow[dr,phantom,"{*}" description]\&
\mathcal{T}_{\rho,\sigma}^{+}\arrow[d]
 \arrow[r,draw=none,"{\deleted{\longrightarrow}}" description]\&\deleted{1}\\
1\arrow[r]\&\mu_2\arrow[r]\arrow[d,equal]\&
R_{\rho,\sigma}\arrow[r,"{\mathrm{can.}}"]\arrow[d,"{\kappa}"]\&
\mathcal{T}_{\rho,\sigma}\arrow[r]\arrow[d,"{\chi}"]\&1\\
1\arrow[r]\&\mu_2\arrow[r]\&G_m\arrow[r,"{\lambda\mapsto\lambda^2}"]\arrow[d]\&
G_m\arrow[r]\arrow[d]\&1\\
\&\&1\&1\&
\end{tikzcd}}

which is nothing but an explicit description of the upper face
\inserted{$(*)$} of the cubical diagrams (96) ---
\deleted{\unsure{so}} \unsure{we} shall take \unsure{note} that this one
\inserted{is \uline{not} cartesian} --- nor the lower face.
Thus, we construct
\deleted{$\unsure{\mathcal{T}_{\rho,\sigma}}$ and
$\unsure{\mathcal{T}_{\rho,\sigma}|R_{\rho,\sigma}}$}

\noindent(107) group $\mathcal{T}_{\rho,\sigma}$, and \inserted{the homomorphisms}
\diagram{\begin{tikzcd}[ampersand replacement=\&,column sep=large]
\mathcal{T}_{\rho,\sigma}\arrow[r,"{\chi}"]\&G_m
\end{tikzcd}},
\diagram{\begin{tikzcd}[ampersand replacement=\&,column sep=large]
\mathcal{T}_{\rho,\sigma}\arrow[r,"{{(\varepsilon_3,\varepsilon_2)}}"]\&\mu_2\times\mu_2
\end{tikzcd}}

[and thereby thanks to (105).--(106)) $R_{\rho,\sigma}$,
$SR_{\rho,\sigma}\simeq\mathcal{T}_{\rho,\sigma}^{+}$, and also
(by \deleted{\unsure{simplifying}}\inserted{replacing}
$\mathcal{T}_{\rho,\sigma}$ by
$\mathcal{T}_{\rho,\sigma}^{\circ}=\operatorname{Ker}
(\mathcal{T}_{\rho,\sigma}\longrightarrow\mu_2\times\mu_2)$)
$R_{\rho,\sigma}^{\circ}$,
$SR_{\rho,\sigma}^{\circ}\simeq\mathcal{T}_{\rho,\sigma}^{+\circ}$
\unsure{$\Rightarrow$}]
in terms of the groups

\noindent(108) $T_\rho,N'_\sigma$ and of
\diagram{\begin{tikzcd}[ampersand replacement=\&,column sep=large]
T_\rho\arrow[r,"{\delta_\rho}"]\&G_m
\end{tikzcd}},
\diagram{\begin{tikzcd}[ampersand replacement=\&,column sep=large,row sep=large]
N'_\sigma\arrow[r,"{\delta_\sigma}"]\&G_m\\
N'_\sigma\arrow[r,"{\varepsilon_\sigma}"]\&\mu_2
\end{tikzcd}},

\mspage{436}

\begin{center}
{\small
436\\
\texttt{https://grothendieck.umontpellier.fr}\\
Université de Montpellier, UMR CNRS 5149, Case courrier 051, Place Eugène Bataillon\\
34095 Montpellier cedex 5 - France
}
\end{center}

\scanpage{438}
\mspage{521}

as a subgroup

\noindent(109)
\diagram{\begin{tikzcd}[ampersand replacement=\&,column sep=small,row sep=large]
\mathcal{T}_{\rho,\sigma}\arrow[r,hook]\&
T_0\arrow[d,"{\delta_0}","{\simeq}"']\&{}\times T_\rho\times N'_\sigma\\
\&G_m\&
\end{tikzcd}}
\[
\begin{aligned}
&\phantom{\{}\,\Vert\\[-2pt]
&\left\{\left(\underset{T_0}{\underline u},
              \underset{T_\rho}{\underline a},
              \underset{N'_\sigma}{\underline b}\right)
 \ \middle|\ 
 \delta_0(\underline u)=\delta_\sigma(\underline b)
 =\varepsilon'\delta_\rho(\underline a),\quad\varepsilon'{}^2=1
 \right\}\\
&\hspace{3em}\simeq
 \left\{\left(\underset{T_\rho}{\underline a},
              \underset{N'_\sigma}{\underline b}\right)
 \ \middle|\
 \delta_\rho(\underline a)=\varepsilon'\delta_\sigma(\underline b),
 \quad\varepsilon'{}^2=1\right\}
\end{aligned}
\]
with
\[
\tag{110}
\left\{
\begin{aligned}
\chi(\deleted{\underline u,}\underline a,\underline b)
 &=\unsure{\delta_\sigma|\delta_0}(\underline b)\\
\varepsilon_3(\deleted{\underline u,}\underline a,\underline b)
 &=\delta_\rho(\underline a)/\delta_\sigma(\underline b)\\
\varepsilon_2(\deleted{\underline u,}\underline a,\underline b)
 &=\varepsilon_\sigma(\underline b)
\end{aligned}
\right.
.
\]

With the help of this, we can construct $\mathcal{M}_{\rho,\sigma}$
by (102), and of course the homomorphism

\noindent(111) $\Theta:$\quad
\diagram{\begin{tikzcd}[ampersand replacement=\&,column sep=large]
\mathcal{M}_{\rho,\sigma}\arrow[r]\&G
\end{tikzcd}}

is deduced by passage to the quotient from the obvious homomorphisms

\noindent(112)
\diagram{\begin{tikzcd}[ampersand replacement=\&,column sep=large]
SG\arrow[r,hook]\&G
\end{tikzcd}},\qquad
\diagram{\begin{tikzcd}[ampersand replacement=\&,column sep=large,row sep=large]
R_{\rho,\sigma}\arrow[rr]\arrow[dr,"{\kappa}"']\&\&G\\
\&G_m\arrow[r,"{\simeq}"']\&
\uline{\operatorname{Centr}}(G)\arrow[u,hook]
\end{tikzcd}}

\deleted{We have not had to}
whence we deduce the composites

\noindent(113)
\diagram{$\left\{\vcenter{\hbox{\begin{tikzcd}[ampersand replacement=\&,column sep=large,row sep=huge]
\mathcal{M}_{\rho,\sigma}\arrow[r,"{\Theta}"]
 \arrow[rr,bend right=25,"{\widetilde{\Theta}}"']\&
 G\arrow[r]\&PG\\
\mathcal{M}_{\rho,\sigma}\arrow[r,"{\Theta}"]
 \arrow[rr,bend right=25,"{\chi}"']\&
 G\arrow[r,"{\det}"]\&G_m
\end{tikzcd}}}\right.$}

We also deduce the exact sequence

\noindent(114)
\diagram{\begin{tikzcd}[ampersand replacement=\&,column sep=large,row sep=large]
1\arrow[r]\&SG\arrow[rr,"{\mathrm{can}}"']\arrow[dr,dashed,"{\mathrm{can}}"']\&\&
\mathcal{M}_{\rho,\sigma}\arrow[r]\&
\mathcal{T}_{\rho,\sigma}\arrow[r]
 \arrow[d,phantom,"{{\rotatebox{90}{$\simeq$}}}" description]\&1\\
\&\&R_{\rho,\sigma}\times SG\arrow[ur,dashed,"{\mathrm{can}}"]\&\&
R_{\rho,\sigma}/\mu_2\&
\end{tikzcd}}

\mspage{437}

\begin{center}
{\small
437\\
\texttt{https://grothendieck.umontpellier.fr}\\
Université de Montpellier, UMR CNRS 5149, Case courrier 051, Place Eugène Bataillon\\
34095 Montpellier cedex 5 - France
}
\end{center}

\scanpage{439}

from the description (102), whence composites

\noindent (115)
\diagram{\begin{tikzcd}[ampersand replacement=\&,column sep=large,row sep=small]
 \& \& T_\rho \\
\mathcal{M}_{\rho,\sigma}\arrow[r]
 \& \mathcal{T}_{\rho,\sigma}
   \arrow[ur,"{\mathrm{pr}_1}"]
   \arrow[dr,"{\mathrm{pr}_2}"]
   \arrow[ddr]
 \& \& \\
 \& \& N'_\sigma\arrow[r,"{\det}"] \& G_m \\
 \& \& \mu_2\times\mu_2
\end{tikzcd}}

In all this, we have not needed the datum \unsure{of|of the} sections $\rho,\sigma$
of $G$, nor to know anything at all about the way in which $T_\rho,N'_\sigma$
\deleted{are}\inserted{have been} obtained from such sections.
But more delicate \unsure{additional} structures of the
$\mathcal{M}_{\rho,\sigma}$ --- such as the datum of the
\unsure{integral} subgroups $\mathcal{M}_{\rho,\sigma}\unsure{(w)|(\omega)}$ ---
depend on them, and will perhaps allow us to recover them.

\uline{\unsure{Remark|Remark}} The cube (96) is the ``composite'' of two cubical
diagrams, which I am going to make explicit, the first of which completes (95) --- which can
be regarded as a partial diagram of this one:

\mspage{438}

\begin{center}
{\small
438\\
\texttt{https://grothendieck.umontpellier.fr}\\
Université de Montpellier, UMR CNRS 5149, Case courrier 051, Place Eugène Bataillon\\
34095 Montpellier cedex 5 - France
}
\end{center}